\chapter{Meta-learning, fast weights, test-time training}
\label{bg:chapter-meta-ttt}

Pre-training과 post-training은 ``배포 전에 무엇을 배울까''를 주로 다룬다. Test-time training과
sleep-time training은 배포 뒤 상태 변화라는 공통점이 있지만 latency path, data window, validation,
persistence가 다르다. 이를 구분하려면 update의 시간척도와 state lifetime을 함께 본다.

\section{Learning to learn: 초기값과 update rule을 학습한다}

\concept{meta-learning}{메타학습}은 여러 task를 경험하며 새 task에 적은 data와 step으로 잘 적응하는
초기값 또는 학습 규칙을 찾는다. MAML의 대표적 형태는 task $\mathcal{T}_i$마다 내부 update

\begin{equation}
\theta'_i=\theta-\alpha\nabla_\theta
\mathcal{L}_{\mathcal{T}_i}^{train}(\theta)
\end{equation}

를 수행한 뒤, 그 결과의 validation loss로 초기값 $\theta$를 갱신한다
\parencite{finn2017maml}.

\concept{bilevel-optimization}{이중수준 최적화}는 ``inner loop에서 adaptation, outer loop에서 그
adaptation 가능성을 평가''하는 구조다.

\begin{equation}
\min_{\theta}\ \mathbb{E}_{\mathcal{T}}
\left[\mathcal{L}^{val}_{\mathcal{T}}
\left(U(\theta,D^{train}_{\mathcal{T}})\right)\right].
\end{equation}

Sleep system에 적용하면 pre/post-training이 좋은 consolidation initializer와 update rule을 배우고,
각 user/device sleep이 local inner loop를 실행할 수 있다. 다만 meta-training task distribution과 실제
user stream이 다르면 ``빠르게 잘못 적응''할 수 있으므로 outer validation에 scope, forgetting, cost를
포함한다.

\section{Fast weights: 상태의 수명을 설계한다}

\concept{fast-weights}{빠른 가중치}는 base weights보다 자주 갱신하고 쉽게 reset하는 parametric state다.
Attention KV cache, recurrent state처럼 gradient 없이 쓰이는 state와 혼동하지 않는다. 여기서는 learned
update로 바뀌며 forward behavior에 들어가는 adapter, memory MLP, normalization parameter 등을 포함한다.

상태를 lifetime으로 나누면 다음과 같다.

\begin{table}[htbp]
\centering
\caption{배포 후 학습 상태의 시간척도}
\label{tab:bg-state-timescale}
\small
\begin{tabularx}{\textwidth}{p{0.18\textwidth} X X X}
\toprule
상태 & update 시점 & persistence & 실패 시 복구 \\
\midrule
query-local & 한 query 내부 & 즉시 폐기 & query 재실행 \\
session fast state & token/session 중 & session 또는 짧은 TTL & session reset \\
user adapter & session 사이 sleep & versioned 장기 & adapter rollback \\
global delta & 여러 cohort 통합 후 & release 단위 & global canary rollback \\
base weights & 드문 offline release & 장기 기준선 & 이전 model release \\
\bottomrule
\end{tabularx}
\end{table}

상태가 오래 살수록 더 많은 validation, provenance, rollback이 필요하다. 빠른 update가 항상 작은
compute라는 뜻도 아니다. token마다 backward하면 작은 parameter라도 latency와 activation memory가
커질 수 있다.

\section{Test-time training의 정확한 경계}

\concept{test-time-training}{테스트 시점 학습}은 test input 자체 또는 가까운 stream에서 자기지도
objective를 만들고 parameter를 갱신해 distribution shift에 적응한다. 초기 연구는 한 unlabeled test
sample의 self-supervised loss로 prediction 전에 update하는 구성을 보였다
\parencite{sun2020ttt}; entropy minimization으로 normalization state를 조정하는 test-time adaptation도
제안되었다 \parencite{wang2021tent}.

다음 세 가지를 구분한다.

\begin{enumerate}
  \item \textbf{TTT-before-answer}: 현재 answer 전에 update가 끝나므로 latency-critical하고 현재 query에
        직접 이익이 있어야 한다.
  \item \textbf{wake consolidation}: session 중 비동기 또는 짧은 idle gap에 최근 stream을 반영한다.
        다음 query부터 유효하지만 아직 wake resource와 충돌한다.
  \item \textbf{sleep training}: 명시적 deferred window에서 더 넓은 replay, validation, compaction,
        publication을 수행한다. 현재 query latency에서 분리된다.
\end{enumerate}

같은 LoRA update도 첫 번째면 TTT, 세 번째면 sleep-time training이다. 따라서 분류 기준은 optimizer나
architecture가 아니라 critical path와 data horizon, persistence, release boundary다.

\section{Reset, accumulation, leakage}

TTT가 sample마다 reset되면 장기기억이 아니고, stream에 누적되면 continual learning이 된다. 누적형은
order sensitivity, adversarial poisoning, user cross-contamination을 받는다. update scope를 query/session/user/
cohort/global로 tag하고 상위 scope로의 승격은 sleep validation을 거치게 한다.

Test input을 학습에 사용한 뒤 같은 sample에서 score를 측정하면 transductive setting이다. 이것을 일반
held-out generalization과 혼동하지 않는다. 더구나 개인 memory에서는 답을 이미 포함한 tool result로
update한 뒤 동일 query를 재평가하면 retrieval cache와 learning gain을 분리할 수 없다.

\begin{keypoint}[Wake와 sleep은 경쟁하는 이름이 아니다]
Wake/test-time update는 즉시 adaptation을, sleep은 replay·compression·cross-session validation과
long-lived publication을 맡는다. 실제 system은 빠른 state가 evidence를 모으고 sleep이 일부를 외부
장기기억이나 parametric artifact로 승격하는 cascade가 된다.
\end{keypoint}
